\booksection{Приложение: Разбор расширения памяти на Android-телефонах}

\chapterintro{При покупке телефона вы часто видите рекламу вроде «12 ГБ + 8 ГБ ОЗУ» — как будто откуда ни возьмись появились дополнительные 8 ГБ оперативной памяти. Но действительно ли это ускоряет телефон? Не вредит ли это ему? В этом разделе мы досконально разберём технологию «расширения памяти».}

\booksubsection{I. Суть определения}

«Расширение памяти» на Android-телефонах (также называется Memory Extension, Memory Fusion, RAM Plus и т.д.)\textbf{не является физическим увеличением чипов оперативной памяти}, а представляет собой системную программную технологию: часть внутреннего пространства памяти (ROM/флеш-память UFS) динамически выделяется и используется как виртуальная оперативная память (RAM).

Приведём аналогию: физическая RAM — это «скоростной рабочий стол», а расширенная память — «временная кладовая». Когда рабочего стола не хватает, система переносит временно неиспользуемые данные в кладовую, а когда нужно — возвращает обратно. Кладовая вместительная, но перенос туда-сюда, конечно, медленнее, чем держать всё прямо на рабочем столе.

\booksubsection{II. Технический принцип}

Основные производители используют два основных способа реализации:

\subsubsection*{Сжатие zRAM}

Внутри RAM выделяется область сжатых данных, куда помещается сжатый код неактивных приложений, степень сжатия примерно 2:1. Данные остаются в RAM, скорость высокая, флеш-память не изнашивается, но требуется участие CPU для сжатия и распаковки.

\subsubsection*{Раздел swap}

Непосредственно во флеш-памяти UFS выделяется фиксированное пространство в качестве виртуальной памяти, куда записываются данные фоновых процессов. Зависит от пропускной способности памяти, скорость чтения/записи намного ниже физической памяти.

Рекламируемая производителями конфигурация «12 ГБ + 8 ГБ ОЗУ» обычно представляет собой комбинацию 12 ГБ физической RAM + 8 ГБ виртуальной памяти на флеш-памяти. Примечательно, что zRAM по умолчанию не занимает встроенную память; только когда zRAM полностью заполнена данными, система создаёт файл подкачки на накопителе, и такая ситуация возникает крайне редко.

\booksubsection{III. Разница скоростей — ключевое узкое место}

Фундаментальное ограничение этой технологии заключается в разнице скоростей. В настоящее время у популярной флеш-памяти UFS 4.0 последовательная скорость чтения составляет около 4200 МБ/с, что составляет лишь 1/10–1/15 от скорости физической памяти LPDDR5X. Поэтому на этапе сжатия zRAM при лёгкой нагрузке скорость близка к реальной RAM, но как только сценарий становится высоконагруженным, скорость резко падает до нескольких десятков раз медленнее.

\booksubsection{IV. Применимые сценарии и границы эффективности}

Расширение памяти не является панацеей, эффект зависит от модели и сценария использования:

\begin{itemize}
  \item \textbf{Полезные случаи}: модели с 6 ГБ физической памяти + UFS 3.1 и выше, свободное пространство $\geq$30 ГБ, при включении расширения на 3 ГБ количество удерживаемых в фоне приложений увеличивается на 3-4, частота перезагрузок снижается на 68\%, подходит для пользователей с интенсивной многозадачностью — курьеров, водителей такси и т.д.
  \item \textbf{Повседневные сценарии}: значительное улучшение для нескольких аккаунтов WeChat, удержания веб-вкладок, ответа на сообщения в мессенджерах — время холодного запуска фоновых приложений заметно сокращается.
  \item \textbf{Ограниченные сценарии}: для загрузки больших игр, редактирования 4K-видео и других сценариев с высокой пропускной способностью улучшение ограничено.
  \item \textbf{Вредные случаи}: модели с 8 ГБ ОЗУ + UFS 2.2 и недостаточным свободным пространством — после включения колебания частоты кадров в играх, наоборот, увеличиваются, просмотр коротких видео расходует больше энергии, переключения тормозят.
  \item \textbf{Бесполезные случаи}: модели с 16 ГБ ОЗУ и более — выгода практически нулевая, максимум ещё одно приложение удерживается в фоне, но при этом нагрузка на CPU увеличивается на 12\%, нагрев растёт на 2-3°C. Для моделей с RAM более 12 ГБ рекомендуется просто отключить функцию во избежание бессмысленных вычислений сжатия.
\end{itemize}

\booksubsection{V. Названия у разных производителей и способ включения}

Разные бренды называют по-разному: Xiaomi — «Расширение памяти», Huawei — «Memory Fusion», Samsung — «RAM Plus», vivo — «Extended RAM», OnePlus — «Memory Extension». Эта функция является встроенной возможностью системы и не имеет прямой привязки к базовой версии Android.

Обычный путь включения: [Настройки] [Расширенные настройки] найти подпункт [Память] или [Производительность] включить «Расширение памяти» и выбрать уровень ёмкости (обычно 2–7 ГБ), вступает в силу после перезагрузки. После включения система будет показывать комбинацию памяти, например «8 ГБ + 4 ГБ».

\begin{tcbwarn}{Примечания}
\begin{itemize}
  \item Расширение памяти занимает часть встроенной памяти, убедитесь, что свободно не менее 20-30 ГБ, иначе это может повлиять на плавность работы системы.
  \item Использование флеш-памяти в качестве виртуальной памяти вызывает износ от чтения/записи; хотя современная флеш-память UFS имеет большой ресурс, длительная высокочастотная подкачка всё же может ускорить старение памяти.
  \item Если физическая память телефона мала (например, 4 ГБ и менее), включение расширения памяти может лишь временно облегчить ситуацию, но не решить проблему тормозов фундаментально.
\end{itemize}

\end{tcbwarn}

\begin{tcbinfo}{Резюме}
Расширение памяти — это зрелая технология системной оптимизации, способная без изменения аппаратного обеспечения повысить возможности многозадачности и удержания фоновых процессов в моделях начального и среднего уровня. Однако она ограничена огромным разрывом в скорости между флеш-памятью и физической памятью и не является «универсальным лекарством». Включать ли её и насколько — нужно решать, исходя из объёма физической памяти вашей модели, спецификации флеш-памяти, свободного места и привычек использования.

\end{tcbinfo}

\begin{tcbpractice}{Практика: проверьте, есть ли на вашем телефоне расширение памяти}
\begin{enumerate}
  \item Откройте [Настройки], найдите «Расширение памяти», «Memory Fusion», «RAM Plus» или «Extended RAM».
  \item Найдя, проверьте, включено ли сейчас, и какой установлен объём расширения.
  \item Сравнивая с содержанием этого раздела, определите: какой объём физической памяти у вашей модели? Какова спецификация флеш-памяти? Достаточно ли свободного места?
  \item Попробуйте отключить или уменьшить объём расширения, понаблюдайте неделю за удержанием фоновых процессов и нагревом, найдите оптимальные для себя настройки.
\end{enumerate}

\end{tcbpractice}
